\usetikzlibrary{patterns,patterns.meta,fadings}
\usetikzlibrary{graphs}
\usetikzlibrary{hobby,backgrounds,calc,trees}
\usetikzlibrary{shapes.misc}
\usetikzlibrary{decorations.pathmorphing,calc,shadows.blur,shadings}

% Font configuration
\usepackage{fontspec}
%\defaultfontfeatures{RawFeature={+axis={wght=100}}}
\setsansfont{Caveat}
\setmainfont{Caveat}
\setmathrm{Caveat}
\setmathsf{Caveat}
\setmathtt{Caveat}

% Pencil line styles

\usepackage{ifthen}
\usetikzlibrary{calc,patterns,decorations,plotmarks}
\pgfdeclaredecoration{penciline}{initial}{
  \state{initial}[width=+\pgfdecoratedinputsegmentremainingdistance, auto corner
    on length=1pt, ]{ \ifthenelse {\pgfkeysvalueof{/tikz/penciline/jag ratio}=0}
    { \pgfpathcurveto%
              {% 1st control point
                \pgfpoint
                    {\pgfdecoratedinputsegmentremainingdistance/2}
                    {2*rnd*\pgfdecorationsegmentamplitude}
              }
              {%% 2nd control point
                \pgfpoint
                %% Make sure random number is always between origin and target points
                    {\pgfdecoratedinputsegmentremainingdistance/2}
                    {2*rnd*\pgfdecorationsegmentamplitude}
              }
              {% 2nd point (1st one is implicit)
                \pgfpointadd
                    {\pgfpointdecoratedinputsegmentlast}
                    {\pgfpoint{0*rand*1pt}{0*rand*1pt}}
              }
        } {
          \pgfpathcurveto%
              {% 1st control point
                \pgfpoint
                    {\pgfdecoratedinputsegmentremainingdistance*rnd*1pt}
                    {\pgfkeysvalueof{/tikz/penciline/jag ratio}*
                      rand*\pgfdecorationsegmentamplitude}
              }
              {%% 2nd control point
                \pgfpoint
                %% Make sure random number is always between origin and target points
                    {(.5+0.25*rand)*\pgfdecoratedinputsegmentremainingdistance}
                    {\pgfkeysvalueof{/tikz/penciline/jag ratio}*
                      rand*\pgfdecorationsegmentamplitude}
              }
              {% 2nd point (1st one is implicit)
                \pgfpointadd
                    {\pgfpointdecoratedinputsegmentlast}
                    {\pgfpoint{rand*1pt}{rand*1pt}}
              }
        }
  }
  \state{final}{}
}

\tikzset{
  penciline/.code={\pgfqkeys{/tikz/penciline}{#1}},
  penciline={
    jag ratio/.initial=5,
    decoration/.initial = penciline,
  },
  penciline/.style = {
    decorate,
    %%decoration={\pgfkeysvalueof{/tikz/penciline/decoration}},
    penciline/.cd,
    #1,
    /tikz/.cd,
  },
  decorate,
  decoration={\pgfkeysvalueof{/tikz/penciline/decoration}},
}

% Crayon fills

\tikzset{
  rect crayon fill/.style args = {#1}{
    append after command={
      \pgfextra
      \draw[Pink,opacity=1.0,decorate,line width=0.9mm,line join=bevel,#1] (\tikzlastnode.north west)
      \foreach \x[remember=\x as \lastx(initially 0)] in {0.1,0.2,...,1.0}{
       -- ($(\tikzlastnode.north west)!\lastx!(\tikzlastnode.south west)$) --($(\tikzlastnode.north west)!\lastx+0.05!(\tikzlastnode.north east)$)
      }
      \foreach \x[remember=\x as \lastx(initially 0)] in {0.1,0.2,...,1}{
       -- ($(\tikzlastnode.south west)!\lastx!(\tikzlastnode.south east)$) --($(\tikzlastnode.north east)!\lastx+0.05!(\tikzlastnode.south east)$)
      };
      \endpgfextra
    }
  }
}

\tikzset{
  axis crayon fill/.style = { insert path = {
      \pgfextra
      \draw[LightBlue!40,opacity=1.0,decorate,line width=0.9mm,line join=bevel] (\tikzlastnode.north west)
      \foreach \x[remember=\x as \lastx(initially 0)] in {0.1,0.2,...,1.0}{
       -- ($(\tikzlastnode.north west)!\lastx!(\tikzlastnode.south west)$) --($(\tikzlastnode.north west)!\lastx+0.05!(\tikzlastnode.north east)$)
      }
      \foreach \x[remember=\x as \lastx(initially 0)] in {0.1,0.2,...,1}{
       -- ($(\tikzlastnode.south west)!\lastx!(\tikzlastnode.south east)$) --($(\tikzlastnode.north east)!\lastx+0.05!(\tikzlastnode.south east)$)
      };
      \endpgfextra
    }
  }
}

\newcommand{\plotcrayonbg}{
    \begin{pgfonlayer}{background}
     \node[rect crayon fill={LightBlue!60, line width=2mm},fit={(current axis)}] {};
    \end{pgfonlayer}
}
% This is the base for the fractal decoration. It takes a random point between the start and end, and
% raises it a random amount, thus transforming a segment into two, connected at that raised point
% This decoration can be applied again to each one of the resulting segments and so on, in a similar
% way of a Koch snowflake.
\pgfdeclaredecoration{irregular fractal line}{init}
{
  \state{init}[width=\pgfdecoratedinputsegmentremainingdistance]
  {
    \pgfpathlineto{\pgfpoint{random*\pgfdecoratedinputsegmentremainingdistance}{(random*\pgfdecorationsegmentamplitude-0.02)*\pgfdecoratedinputsegmentremainingdistance}}
    \pgfpathlineto{\pgfpoint{\pgfdecoratedinputsegmentremainingdistance}{0pt}}
  }
}


% define some styles
\tikzset{
   paper/.style={draw=Orange!60,
                 fill=Orange!60,
                 % blur shadow,
                 % lower left=black!20, upper left=black!15, upper right=white, lower right=black!10
                 },
   irregular border/.style={decoration={irregular fractal line, amplitude=0.2},
           decorate,
     },
   ragged border/.style={ decoration={random steps, segment length=7mm, amplitude=2mm},
           decorate,
   }
}

% Macro to draw the shape behind the text, when it fits completly in the
% page
\def\tornpaper#1{
\tikz{
  \node[inner sep=1em] (A) {#1};  % Draw the text of the node
  \begin{pgfonlayer}{background}  % Draw the shape behind
  \fill[paper] % recursively decorate the bottom border
        decorate[irregular border]{decorate{decorate{decorate{decorate[ragged border]{
        ($(A.south east) $) -- ($(A.south west)$)
        -- (A.north west) -- (A.north east) -- cycle
        }}}}};
  \end{pgfonlayer}}
}

\tikzset{
  ripped lit/.style args = {#1}{
    append after command={
      \pgfextra
      \fill[paper,#1] % recursively decorate the bottom border
            decorate[irregular border]{decorate{decorate{decorate{decorate[ragged border]{
            ($(\tikzlastnode.south east) $) -- ($(\tikzlastnode.south west)$)
            -- (\tikzlastnode.north west) -- (\tikzlastnode.north east) -- cycle
            }}}}};
      \endpgfextra
    }
  }
}
